% ============================================
% Survey Structure Figure Template
% 使用Forest包生成多层级彩色分类树
% ============================================

% Define colors and styles
\definecolor{hidden-draw}{RGB}{128, 128, 128}
<define_color>

% 大纲总样式
\tikzstyle{ot}=[
    rectangle,
    draw=hidden-draw,
    rounded corners,
    align=left,
    text opacity=1,
    text=black,
    minimum height=2.5em,
    minimum width=5em,
    inner xsep=4pt,
    inner ysep=4pt,
    line width=0.8pt,
    font=\small,
]

% 大纲主样式-root
\tikzstyle{root}=[
    ot,
    draw=gray!80,
    fill=gray!35,
    font=\small\bfseries,
]

% 大纲主样式-1级 (Section)
\tikzstyle{ot1}=[
    ot,
    minimum height=2em,
    font=\small,
]

% 大纲主样式-2级 (Subsection)
\tikzstyle{ot2}=[
    ot,
    minimum height=1.8em,
    font=\footnotesize,
]

% 大纲主样式-3级 (Subsubsection/Leaf with refs)
\tikzstyle{ot3}=[
    ot, 
    minimum height=1.5em,
    font=\footnotesize,
]

% 引用样式 (用于显示cite角标)
\tikzstyle{refnode}=[
    rectangle,
    draw=gray!50,
    rounded corners=2pt,
    fill=gray!10,
    font=\scriptsize,
    inner xsep=2pt,
    inner ysep=1pt,
    text=gray!70,
]

\begin{figure*}[!th]
    \centering
    \resizebox{0.98\textwidth}{!}
    {
        \begin{forest}
            for tree={
                grow=east,
                reversed=true,
                anchor=base west,
                parent anchor=east,
                child anchor=west,
                base=left,
                font=\small,
                rectangle,
                draw=hidden-draw,
                rounded corners,
                align=left,
                minimum width=2em,
                edge+={darkgray, line width=1pt},
                s sep=7pt,
                inner xsep=4pt,
                inner ysep=3pt,
                line width=0.8pt,
                ver/.style={rotate=90, child anchor=north, parent anchor=south, anchor=center},
            }, 
            <tree_content>
        \end{forest}
    }
    \caption{<caption>}
    \label{fig:<label>}
\vspace{-0.3cm}
\end{figure*}
